\documentclass[12pt]{article}
\usepackage[a4paper,margin=1in]{geometry}
\usepackage[utf8]{inputenc}
\usepackage[T1]{fontenc}
\usepackage{amsmath,amssymb}
\usepackage{textcomp}
\usepackage{parskip}
\begin{document}
\begin{center}
{\Large \textbf{2014 STEP I}}
\end{center}

\subsection*{Page 1}
\begin{flushleft}
\bigskip\textbf{Section A: Pure Mathematics}\par
1 All numbers referred to in this question are non-negative integers.\par
\textbf{(i)} Express each of the numbers 3, 5, 8, 12 and 16 as the difference of two non-zero squares.\par
\textbf{(ii)} Prove that any odd number can be written as the difference of two squares.\par
\textbf{(iii)} Prove that all numbers of the form 4k, where k is a non-negative integer, can be written\par
as the difference of two squares.\par
\textbf{(iv)} Prove that no number of the form 4k + 2, where k is a non-negative integer, can be\par
written as the difference of two squares.\par
\textbf{(v)} Prove that any number of the form pq, where p and q are prime numbers greater than 2,\par
can be written as the difference of two squares in exactly two distinct ways. Does this\par
\[result hold if p is a prime greater than 2 and q = 2?\]
\textbf{(vi)} Determine the number of distinct ways in which 675 can be written as the difference\par
of two squares.\par
\bigskip\textbf{Question 2.} (i) Show that ln(2 x) dx = (2 x) ln(2 x) + (2 x) + c , where x < 2.\par
- - - - -\par
\[\leq\]
\textbf{(ii)} Sketch the curve A given by y = ln x2 4 .\par
| - |\par
\textbf{(iii)} Show that the area of the finite region enclosed by the positive x-axis, the y-axis and\par
\[the curve A is 4 ln(2 + \sqrt\{3\}) 2\sqrt\{3\} .\]
-\par
\textbf{(iv)} The curve B is given by y = ln x2 4 . Find the area between the curve B and the\par
| - |\par
\[x-axis with x < 2. \geq \geq\]
\[| | \geq \geq\]
[Note: you may assume that t ln t 0 as t 0.]\par
\to \to\par
\end{flushleft}
\subsection*{Page 2}
\begin{flushleft}
\bigskip\textbf{Question 3.} The numbers a and b, where b > a \leq 0, are such that\par
b b 2\par
\[x dx = x dx .\]
\[(a )(a \int\]
\textbf{(i)} In the case a = 0 and b > 0, find the value of b.\par
\textbf{(ii)} In the case a = 1, show that b satisfies\par
3 2\par
\[3b b 7b 7 = 0 .\]
- - -\par
Show further, with the help of a sketch, that there is only one (real) value of b that\par
satisfies this equation and that it lies between 2 and 3.\par
\textbf{(iii)} Show that 3p2 + q2 = 3p2q, where p = b + a and q = b a, and express p2 in terms\par
\[of q. Deduce that 1 < b a \geq 4 . -\]
- 3\par
\bigskip\textbf{Question 4.} An accurate clock has an hour hand of length a and a minute hand of length b (where b > a),\par
both measured from the pivot at the centre of the clock face. Let x be the distance between\par
\[the ends of the hands when the angle between the hands is θ, where 0 \geq θ < \pi.\]
\[Show that the rate of increase of x is greatest when x = (b^2 a^2) 2 1 .\]
-\par
In the case when b = 2a and the clock starts at mid-day (with both hands pointing vertically\par
upwards), show that this occurs for the first time a little less than 11 minutes later.\par
\bigskip\textbf{Question 5.} (i) Let f(x) = (x + 2a)3 27a2x, where a \leq 0. By sketching f(x), show that f(x) \leq 0\par
-\par
\[for x \leq 0.\]
\textbf{(ii)} Use part (i) to find the greatest value of xy2 in the region of the x-y plane given by\par
x \leq 0, y \leq 0 and x+2y \geq 3 . For what values of x and y is this greatest value achieved?\par
\textbf{(iii)} Use part (i) to show that (p + q + r)3 \leq 27pqr for any non-negative numbers p, q and r.\par
\[If (p + q + r)3 = 27pqr, what relationship must p, q and r satisfy?\]
\end{flushleft}
\subsection*{Page 3}
\begin{flushleft}
\bigskip\textbf{Question 6.} (i) The sequence of numbers u0, u1, . . . is given by u0 = u and, for n \leq 0,\par
\[un+1 = 4un(1 un) . ( )\]
- ∗\par
\[In the case u = sin θ for some given angle θ, write down and simplify expressions for\]
\[u^1 and u^2 in terms of θ. Conjecture an expression for un and prove your conjecture.\]
\textbf{(ii)} The sequence of numbers v0, v1, . . . is given by v0 = v and, for n \leq 0,\par
\[vn+1 = pvn + qvn + r ,\]
-\par
\[where p, q and r are given numbers, with p = 0. Show that a substitution of the form\]
)\par
\[vn = αun + β, where α and β are suitably chosen, results in the sequence ( ) provided\]
∗\par
that\par
\[4pr = 8 + 2q q .\]
-\par
\[Hence obtain the sequence satisfying v^0 = 1 and, for n \leq 0, vn+1 = vn 2 + 2vn + 2 .\]
-\par
\bigskip\textbf{Question 7.} In the triangle OAB, the point D divides the side BO in the ratio r : 1 (so that BD = rDO),\par
\[and the point E divides the side OA in the ratio s : 1 (so that OE = sEA), where r and s\]
are both positive.\par
\textbf{(i)} The lines AD and BE intersect at G. Show that\par
rs 1\par
\[g = a + b ,\]
1 + r + rs 1 + r + rs\par
where a, b and g are the position vectors with respect to O of A, B and G, respectively.\par
\textbf{(ii)} The line through G and O meets AB at F . Given that F divides AB in the ratio t : 1,\par
find an expression for t in terms of r and s.\par
\bigskip\textbf{Question 8.} Let La denote the line joining the points (a, 0) and (0, 1 a), where 0 < a < 1. The line Lb\par
-\par
is defined similarly.\par
\textbf{(i)} Determine the point of intersection of La and Lb, where a = b.\par
)\par
\textbf{(ii)} Show that this point of intersection, in the limit as b a, lies on the curve C given by\par
\to\par
\[y = (1 \sqrt\{x\}) 2 (0 < x < 1) .\]
-\par
\textbf{(iii)} Show that every tangent to C is of the form La for some a.\par
\end{flushleft}
\subsection*{Page 4}
\begin{flushleft}
\bigskip\textbf{Section B: Mechanics}\par
\bigskip\textbf{Question 9.} A particle of mass m is projected due east at speed U from a point on horizontal ground\par
\textasciicircum{}\circ\par
\[at an angle θ above the horizontal, where 0 < θ < 90 . In addition to the gravitational\]
force mg, it experiences a horizontal force of magnitude mkg, where k is a positive constant,\par
acting due west in the plane of motion of the particle. Determine expressions in terms of U ,\par
θ and g for the time, TH , at which the particle reaches its greatest height and the time, TL,\par
at which it lands.\par
Let T = U cos θ/(kg). By considering the relative magnitudes of TH , TL and T , or otherwise,\par
1 1\par
sketch the trajectory of the particle in the cases k tan θ < , < k tan θ < 1, and k tan θ > 1.\par
2 2\par
\[What happens when k tan θ = 1?\]
\bigskip\textbf{Question 10.} (i) A uniform spherical ball of mass M and radius R is released from rest with its centre\par
a distance H + R above horizontal ground. The coefficient of restitution between the\par
ball and the ground is e. Show that, after bouncing, the centre of the ball reaches a\par
\[height R + He2 above the ground.\]
\textbf{(ii)} A second uniform spherical ball, of mass m and radius r, is now released from rest\par
together with the first ball (whose centre is again a distance H + R above the ground\par
when it is released). The two balls are initially one on top of the other, with the second\par
ball (of mass m) above the first. The two balls separate slightly during their fall, with\par
their centres remaining in the same vertical line, so that they collide immediately after\par
the first ball has bounced on the ground. The coefficient of restitution between the\par
balls is also e. The centre of the second ball attains a height h above the ground.\par
\[Given that R = 0.2, r = 0.05, H = 1.8, h = 4.5 and e = , determine the value\]
of M/m.\par
\end{flushleft}
\subsection*{Page 5}
\begin{flushleft}
\bigskip\textbf{Question 11.} The diagrams below show two separate systems of particles, strings and pulleys. In both\par
systems, the pulleys are smooth and light, the strings are light and inextensible, the particles\par
move vertically and the pulleys labelled with P are fixed. The masses of the particles are as\par
indicated on the diagrams.\par
P P\par
P\par
M\par
m m\par
M m\par
System I System II\par
\textbf{(i)} For system I show that the acceleration, a1, of the particle of mass M , measured in\par
the downwards direction, is given by\par
M m\par
\[a^1 = - g ,\]
M + m\par
where g is the acceleration due to gravity. Give an expression for the force on the\par
pulley due to the tension in the string.\par
\textbf{(ii)} For system II show that the acceleration, a2, of the particle of mass M , measured in\par
the downwards direction, is given by\par
M 4μ\par
\[a^2 = - g ,\]
M + 4μ\par
\[m1m2\]
\[where μ = .\]
\[m^1 + m^2\]
\[In the case m = m^1 + m^2, show that a^1 = a^2 if and only if m^1 = m^2.\]
\end{flushleft}
\subsection*{Page 6}
\begin{flushleft}
\bigskip\textbf{Section C: Probability and Statistics}\par
\bigskip\textbf{Question 12.} A game in a casino is played with a fair coin and an unbiased cubical die whose faces are\par
labelled 1, 1, 1, 2, 2 and 3. In each round of the game, the die is rolled once and the coin\par
is tossed once. The outcome of the round is a random variable X. The value, x, of X is\par
determined as follows. If the result of the toss is heads then x = ks 1 , and if the result of\par
| - |\par
\[the toss is tails then x = k s , where s is the number on the die and k is a given number.\]
\[Show that E(X^2) = k + 1 | 3(k - | 1)2/6.\]
-\par
Given that both E(X2) and E(X) are positive integers, and that k is a single-digit positive\par
integer, determine the value of k, and write down the probability distribution of X.\par
A gambler pays \pounds 1 to play the game, which consists of two rounds. The gambler is paid:\par
\pounds  w, where w is an integer, if the sum of the outcomes of the two rounds exceeds 25;\par
\pounds 1 if the sum of the outcomes equals 25;\par
nothing if the sum of the outcomes is less that 25.\par
Find, in terms of w, an expression for the amount the gambler expects to be paid in a game,\par
and deduce the maximum possible value of w, given that the casino's owners choose w so\par
that the game is in their favour.\par
\bigskip\textbf{Question 13.} A continuous random variable X has a triangular distribution, which means that it has a\par
probability density function of the form\par
\[g(x) for a < x \geq c\]
\[f(x) = ⎧ h(x) for c \geq x < b\]
⎪\par
⎨\par
0 otherwise,\par
⎪\par
⎩\par
where g(x) is an increasing linear function with g(a) = 0, h(x) is a decreasing linear function\par
\[with h(b) = 0, and g(c) = h(c).\]
2(x a)\par
\[Show that g(x) = - and find a similar expression for h(x).\]
(b a)(c a)\par
- -\par
\textbf{(i)} Show that the mean of the distribution is (a + b + c).\par
\textbf{(ii)} Find the median of the distribution in the different cases that arise.\par
\end{flushleft}
\end{document}
